\subsection{The integral divided norm and the two remaining orbit-twelve compact blocks}

\begin{statement}{D.65}{Proved}{An actual source correction fills two further arithmetic compact blocks}
Keep the marked sources and all comparison maps of D.64:
\[
 X=\widehat B/(R_{50}\cap L_2),\qquad Y=X/\mathcal K,
 \qquad X_{64}=X/\mathcal J_{64},\qquad
 \mathcal L=\ker(U^-\longrightarrow U_*^-).
 \tag{D.65.1}
\]
In particular, \(G\) is the actual arithmetic group of order
\(576\cdot3^{97}\), and \(t=sr\) has order 36 in \(G\).
The maps to
\[
 G^\dagger=(\mathbb Z/3)^4\rtimes Q\longrightarrow G_*\qquad
 (|G^\dagger|=46656, |G_*|=3888)
 \tag{D.65.2}
\]
are the actual D.60 maps. Use the integral compact lattices \(M_a\),
their small maps \(\phi_a\), and the central modules \(C_a\) of D.64.
Only \(a=2,3\), the two orbits of size twelve, enter the new correction.

There is an actual continuous peripheral automorphism of
\[
 \begin{aligned}
 X_{65}&=X/\mathcal J_{65},\\
 \mathcal J_{65}
   &=\ker\left(\mathcal J_{63}
       \xrightarrow{(q_1,q_2,q_3,\nu_4q_4)}
                  C_1\oplus C_2\oplus C_3\oplus\nu_4(C_4)\right).
 \end{aligned}
 \tag{D.65.3}
\]
It has multiplier one, fixes the original cusp \(t\), is identity on
\(G\), and retains the complete D.57 action with its literal small
conjugators. Its \(\mathcal W_h\) action, normalized parameter, and all
proved metabelian and CM comparisons through that map are unchanged.
Both actual arithmetic branches and all their retained framings are
included.

Here are the exact new maps and dimensions. The two D.64 retractions
are denoted by \(p_a:C_a\to I_a\), with \(\dim I_a=4\). Then
\[
 \begin{aligned}
 \mathcal J_{65}&\subset\mathcal J_{64},\\
 \ker(X_{65}\to X_{64})
       &=\mathcal J_{64}/\mathcal J_{65}
          \simeq\ker p_2\oplus\ker p_3,\\
 \dim\ker(X_{65}\to X_{64})&=5752,\qquad
 \dim\bigl(\ker(X_{65}\to X_{64})\bigr)_G=12,\\
 \dim(\mathcal J_{63}/\mathcal J_{65})&=7682,
       \qquad\dim(\mathcal J_{63}/\mathcal J_{65})_G=26.
 \end{aligned}
 \tag{D.65.4}
\]
The action on the lower quotient is the new action described below;
its comparison with the chosen D.64 action is retained explicitly.

Let \(\mathcal O=\mathbb Z_2[\omega]\),
\(\omega^2+\omega+1=0\). On the actual rank-16
\(\mathcal O\)-lattice \(M_a\), write \(H_a\) for the unordered pair
stabilizer, of order 48. In integral Hermitian coordinates its exterior
action is
\[
 \rho_a(g)B=
 \begin{cases}
 A_gB\bar A_g^{\,t},&g\text{ preserves the character},\\
 -A_g\bar B\bar A_g^{\,t},&g\text{ exchanges the character}.
 \end{cases}
 \tag{D.65.5}
\]
Both the whole integral exterior norm and the whole integral
endomorphism trace are even. Consequently the maps
\[
 \begin{aligned}
 n_a(\bar B)&=\frac12\sum_{g\in H_a}\rho_a(g)B\pmod2,\\
 k_a(\bar A)&=\frac12\sum_{g\in H_a}gAg^{-1}\pmod2,
 \qquad A:M_a\longrightarrow\ker\phi_a
 \end{aligned}
 \tag{D.65.6}
\]
are independent of the indicated integral lifts. The first map descends
to the actual coinvariants of the D.64 kernel. Its exact ranks are
\[
 \begin{array}{c|r|r|r|r}
 a&\dim(K_a)_{H_a}&\operatorname{rank}n_a
      &\dim\ker n_a&\operatorname{rank}k_a\\\hline
 2&10&6&4&16\\
 3&10&6&4&16
 \end{array}
 \tag{D.65.7}
\]
In each row \(\ker n_a\) is exactly the D.64 actual source-involution
image. For every permitted finite framing, the actual compact reduction
\(D_a\) is explicitly determined and is an isometry. With
\(\Omega_a=\overline{(H_a^{\rm int})^{-1}}\), the map
\[
 \bar A\longmapsto
 k_a(\bar A)\Omega_a\bar D_a^{\,t}
             +D_a\Omega_a\overline{k_a(\bar A)}^{\,t}
 \tag{D.65.8}
\]
has precisely the rank-six image of \(n_a\). Complete coefficient
matrices and right inverses on all ten coinvariant representatives are
retained in the certificate.
Here \(H_a^{\rm int}\) denotes the D.64 intersection matrix after
reduction modulo two; \(\Omega^{\rm int}\) in the integral boundary
equation below is the bivector before reduction.

The integral homology changes realizing (D.65.8) factor through actual
finite free \(\mathbb Z_2[G]\)-lattices. They have cusp-fixed peripheral
derivations \(\eta_a:Y\to\mathcal L\). Put
\[
 \beta_Y(g)=(\eta_2(g)+\eta_3(g))\alpha_{64,Y}(g).
 \tag{D.65.9}
\]
The coefficients \(\eta_a\) are selected by the retained right inverse
applied to the actual divided-norm cusp error. The proof constructs
their two literal source conjugators, at all integral homology levels
simultaneously. After the earlier self-block repairs, the resulting
\(Y\)-action is exactly
\[
 \operatorname{Int}(b_{t,Y})\beta_Y,
 \qquad b_t=\prod_j\prod_{i=0}^{35}t^iu_jt^{-i},
 \qquad u_j\in S_{57}.
 \tag{D.65.10}
\]
Here only the transported earlier extra-projective self corrections
enter \(b_t\). Its \(Y\)-image centralizes \(t\), has zero small and
\(\mathcal W_h\) images, and acts trivially on integral homology.
The subsequent eight new source-involution corrections and central
augmentation corrections leave this \(Y\)-action unchanged.

There is also an exact limit on the larger fixed-homology centralizer
route. On the actual lattice \(\mathcal L\),
\[
 \mathcal L^s=(1+s)\mathcal L,\qquad
 \mathcal L^t=N_{36}(t)\mathcal L,\qquad
 \ker N_{36}(t)=(1-t)\mathcal L.
 \tag{D.65.11}
\]
Thus all homology centralizer changes in the kernel of the literal
small comparison are represented by the norm words already tested in
D.64. The new lift uses the actual homology change (D.65.9).

The remaining source is explicit:
\[
 \begin{aligned}
 \ker(X\to X_{65})&=\mathcal J_{65},\qquad
 \ker(X_{65}\to Y)=\mathcal K/\mathcal J_{65},\\
 \dim\mathcal J_{65}&=\dim\mathcal J_{64}-5752,\\
 \dim(\mathcal J_{65})_{\text{specified small characters}}&=276478,\\
 S_{65}&=\{b\in R_{50}:
             \operatorname{im}_{\mathcal T_2}(b)\in\mathcal J_{65}\},\qquad
 S_{65}/(R_{50}\cap L_2)=\mathcal J_{65}.
 \end{aligned}
 \tag{D.65.12}
\]
It still maps onto \(\ker\nu_4\), of dimension 1918, and is nonzero.
The two remaining old-orbit compact norm-kernel coordinates survive
the present test. This is a statement about the specified compact
test, not a formula for \((\mathcal J_{65})_G\). The boundary in the
remaining source kernel and the full profinite realization remain
unresolved. The manuscript passages are Part D(e), section 51.16,
folio 734 verso, scan 174, and section 51.15, folios 725--727,
scans 156--160.
\end{statement}
\begin{proof}
\emph{The actual compact action and its finite framing.}
All coefficient modules start from the marked graph in (D.65.2).
The original \(s,r\) edges, the affine \(r\)-translation
\((1,1,1,0)\), the recorded free edges, and the original character
representatives are retained. On a new character fibre the affine,
cusp and compact ranks are \(96,48,48\) over \(\mathbb F_4\).
The compact decomposition has the original rank-32 projective part
and the original rank-16 complement.

The two possible D.57 branches have \(\beta\)-images in \(G_*\)
with zero translation and linear states 26 and 33. Their literal
centralizing lifts in \(Q\) are respectively
\[
             \{476,498\},\qquad\{456,461\}.
 \tag{D.65.13}
\]
The \(r\)-conjugator has image one in \(G\), by the retained odd-prime
identity component. All later homology corrections lie above that
same finite group. Thus its image in \(G^\dagger\) is still one.
The translation coordinate of \(\beta\) in \(G^\dagger\) is zero:
the map to \(G_*\) is identity on this translation subgroup. These
facts give the complete four-element list (D.65.13), including any
remaining finite centralizer ambiguity.

The program \path{d_compact_rigid_source.py} constructs an actual
equivariant affine section of the rank-16 compact complement. On
the unordered Sylow subgroup of order 16, the cusp lattice is free
of rank six over its binary group ring. Six literal regular
generators and all their 96 real orbit columns are retained. Extend
their coefficient functionals to the affine lattice, coinduce,
average over the three Sylow cosets, and average over the odd
coefficient-scalar group. The resulting cusp retraction is checked
on every column. Subtracting it gives the asserted affine section;
all 48 semilinear intertwiners are checked.

For a literal word \(\beta\) with a prescribed image in (D.65.13),
the actual edge map is right convolution by
\[
                           e_1-p(\beta).
 \tag{D.65.14}
\]
Indeed the general source-conjugator edge map is
\(p(A)-p(\beta)+A_Ge_1\), and here \(A_G=1\).
Every path term retains its four-dimensional translation phase.
The program finds an actual word by breadth-first search on all
46656 marked states, evaluates its full path, applies (D.65.14) to
the actual affine section, and then compactifies. All eight results
have rank 16, satisfy the marked small square, commute with all
semilinear stabilizer maps, and preserve the entire D.64 form.

This determines the compact action of the actual source, including
its possible higher-kernel conjugator changes. A difference between
two edge maps with the same finite conjugator coordinates factors
through the regular edge lattice. Its compact \(M_a\)-to-\(M_a\)
component therefore factors through a projective module. The Higman
trace description follows directly by writing the two maps into
and out of a regular module in its group basis. On these actual
modules every one of the 512 real coefficient columns of that full
trace is zero. This is checked in \path{d_compact_full_trace.py}.
Odd-character averaging first reduces a general real map to the
recorded coefficient-linear map; the full stabilizer trace is then
exactly the checked one. Inflation across the odd kernel of
\(G\to G^\dagger\) introduces only an invertible odd norm factor.
Thus every such difference is zero on the compact reduction.
The computed matrices are the actual \(D_a\).

There is a further exact module comparison. Let \(H_a^0\simeq S_4\)
be the ordered stabilizer and let \(J_a\simeq S_3\) be the subgroup
generated by the retained order-three element and its recorded
normalizing involution. The fixed subspace of \(\overline M_a\)
has rank eight and maps onto the small rank-four module. A literal
section of that map, followed by trace-one semilinear averaging,
gives
\[
 \mathbb F_4[H_a^0/J_a]\otimes\overline P_{*,a}
                 \xrightarrow{\ \sim\ }\overline M_a.
 \tag{D.65.15}
\]
The certificate retains the section, all four coset columns, the
inverse matrix, every stabilizer action, and both diagonal and
off-diagonal blocks of the actual pairing. The small map is the
sum of the four coordinates. This proves the comparison by its
maps, and also gives a direct structural check on the zero trace.

\emph{The integral graph and the signed exterior lattice.}
The program \path{d_compact_integral.py} works in
\(\mathcal O/4\), with the literal encoding
\(a+4b\) for \(a+b\omega\). It starts from the actual D.64 integral
cycle lifts. The chosen free-edge rows are identity, so a pure
\(Q\)-action is obtained by its actual edge permutation and, on an
exchanging element, coefficient conjugation. Every resulting cycle
identity and every reduction to the old marked matrix is checked.
The comparison with the full marked graph of \(G^\dagger\) is the
literal Fourier map
\(\sum_v a_{q,v}[q,v]\mapsto
\sum_v\omega^{\ell v}a_{q,v}[q]\). On the selected real character
pair its inverse has coefficient
\(3^{-4}\operatorname{Tr}_{\mathcal O/\mathbb Z_2}
(\omega^{-\ell v}u_q)\). Orthogonality of the two nontrivial
characters proves both composites, and the actual edge phases make
this a chain map commuting with the cusp maps. Thus the twisted
graph carries the integral marked module, with its specified
coefficient and deck framing.

The cusp columns use the signed word \((sr)^{12}\) in
\(G^\dagger\), with all twelve translations. A unit-pivot quotient
gives the integral compact lattice. The four projective generators
are lifted with the order-three eigenprojector and the integral
trace-one coefficient \(-\omega\), whose sum with its conjugate is
one. Coinduction of the lifted eigenfunctionals gives a retraction;
correct its composite by its inverse and then apply the same
trace-one exchange average. Its reduction is the original D.61
retraction. The complementary columns give the actual integral
rank-16 action matrices. All 144 generating group products per
orbit are checked modulo four.

These are reductions of actual integral lattices with the stated
marking. Projectivity lifts the included projective lattice and
its prescribed modulo-four inclusion. Relative injectivity of its
dual lifts the prescribed retraction, exactly as in D.61. Correcting
the composite by its inverse retains the given reduction modulo
four. If two such integral projectors \(e,e'\) are chosen, then
\[
             u=e'e+(1-e')(1-e)
 \tag{D.65.16}
\]
is an equivariant isomorphism, congruent to one modulo two, carrying
their two decompositions into each other. Both projective images
have zero small image, so the small map satisfies \(\phi u=\phi\).
Hence these choices are compared by an actual isomorphism preserving
the small marking. Conjugating by a matrix congruent to one modulo
two does not change either divided map in (D.65.6): its additional
terms contain the whole trace or norm, which is already even.

For completeness, the integral Hermitian realization is checked
independently on the underlying real exterior lattice. Put
\(u_i=1\), \(v_i=\omega\) in the \(i\)-th coefficient copy. The
primitive embedding sends a diagonal coordinate to \(u_i\wedge v_i\)
and the off-diagonal coordinates to
\[
 \begin{aligned}
 1&\longmapsto u_i\wedge v_j-v_i\wedge u_j,\\
 \omega&\longmapsto u_i\wedge u_j+v_i\wedge u_j+v_i\wedge v_j.
 \end{aligned}
 \tag{D.65.17}
\]
The coefficients of \(u_i\wedge v_i\), \(u_i\wedge v_j\), and
\(u_i\wedge u_j\) give its integral left inverse. Multiplication
by \(a+b\omega\) has real matrix
\(\left(\begin{smallmatrix}a&-b\\b&a-b\end{smallmatrix}\right)\),
whereas coefficient conjugation has matrix
\(\left(\begin{smallmatrix}1&-1\\0&-1\end{smallmatrix}\right)\).
Their exterior actions give exactly (D.65.5), including its minus
sign. The independent program \path{d_compact_integral_exterior.py}
checks all 12288 basis images per orbit, with exact encoding and
decoding in the ambient rank-496 exterior lattice.

The whole exterior norm is even, as is the whole endomorphism trace.
Replacing an integral lift by twice another vector therefore changes
the divided reduction by zero. This proves the lift independence in
(D.65.6), including for maps to the saturated kernel of \(\phi_a\).
The norm also annihilates every integral augmentation relation, so
its divided reduction descends to the claimed coinvariants.

\emph{The complete divided correction calculation.}
The full divided exterior norm has rank twelve on each rank-256
Hermitian module. On the actual rank-240 kernel it has rank six.
All 11520 kernel action columns per orbit are checked for invariance.
The ten original coinvariant representatives are retained. Each of
the 288 actual source-involution correction columns is killed by
the divided norm. Their rank is four, proving exact equality with
its kernel.

The full divided endomorphism trace has rank 32. Restricting its
input to \(\operatorname{Hom}_{\mathbb F_4}
(\overline M_a,\ker\bar\phi_a)\) uses all 384 real coefficient
columns and gives rank 16. Its outputs still have zero small image.
For every actual matrix \(D_a\), the complete linearized bivector
columns in (D.65.8) span the same six-dimensional subspace as the
divided norm. A right inverse is recorded for each of all ten
coinvariant representatives and checked by exact multiplication.
Thus the certificate supplies an actual linear section on the
entire image, rather than a search for one favorable boundary.

\emph{The actual integral source extension.}
We now realize these coefficient maps in the original source.
For an ordinary integral coefficient-linear lift
\(A:M_a\to\ker\phi_a\), extend it by zero to the other character
pairs and compact summands. The surjection in D.64.11 lifts its
values into \(\mathcal L\). Odd-character projection retains the
selected orbit. Denote its actual global compact summand by
\(\mathsf M_a\). Coinduction gives a factorization
\[
 U\xrightarrow{\pi_{M_a}}\mathsf M_a
       \xrightarrow{j}P\xrightarrow{q}\mathcal L,
 \tag{D.65.18}
\]
where \(P\) is a finite free \(\mathbb Z_2[G]\)-lattice.
Explicitly one may use
\(P=\mathbb Z_2[G]\otimes_{\mathbb Z_2}\mathsf M_a\), with \(G\)
acting only on its group-ring factor, and
\(j(v)=\sum_g g\otimes g^{-1}v\). Define \(q\) by the chosen
ordinary lift on the identity summand and induce it equivariantly.
Multiplying \(q\) by \(3^{-97}\in\mathbb Z_2\) removes the normal
odd-group multiplicity. On the original character fibre the compact
composite is exactly \(\operatorname{Tr}_{H_a}(A)\). The first map
in (D.65.18) kills every cusp class and all other compact summands;
the compact image of the last map is in the selected \(M_a\)-kernel.
Write \(j_U=j\pi_{M_a}:U\to P\) for that first composite.

Write \(E=\mathbb Z_2[G]\) for the actual Bass--Serre edge lattice.
The marked chain sequence is
\[
 0\longrightarrow U\longrightarrow E
   \longrightarrow\mathbb Z_2[G/\langle s\rangle]
          \oplus\mathbb Z_2[G/\langle r\rangle]
   \longrightarrow\mathbb Z_2\longrightarrow0.
 \tag{D.65.19}
\]
The quotient of \(E\) by \(U\) is a free \(\mathbb Z_2\)-lattice.
A regular group-ring lattice is coinduced; the coefficient-of-one
isomorphism identifies maps to it with ordinary lattice
functionals. Extend those functionals across this split lattice
sequence and coinduce back. This proves that the first composite
\(j_U:U\to P\) in (D.65.18) extends to a map \(\ell:E\to P\).
It is the required relative injectivity assertion with its actual
extension map, not an assertion that all integral modules are
projective.

The actual path cocycle has
\(p(s)=0\), \(p(r)=e_1-e_r\). Put \(w=\ell(e_1)\).
Since (D.65.18) kills the full cusp homology,
\[
       N_{36}(t)\ell(p(t))
          =\ell(p(t^{36}))=j_U([t^{36}])=0.
 \tag{D.65.20}
\]
The restriction of \(P\) to \(\langle t\rangle\) is free. In each
regular cyclic orbit, the sum-zero coefficients are exactly the
successive differences. This gives \(b\in P\) with
\(\ell(p(t))=(t-1)b\). Therefore
\[
 \begin{aligned}
 \eta_P(g)&=\ell(p(g))+(1-g)b,\\
 \eta_P(t)&=0,\qquad
 \eta_P(s)=(1-s)b,\qquad
 \eta_P(r)=(1-r)(w+b),\\
 \eta_P(n)&=j_U([n])\quad(n\in\ker(\widehat B\to G)).
 \end{aligned}
 \tag{D.65.21}
\]
Composition with \(q\) gives the actual \(\eta_a\).
Its restriction to any defining kernel of \(Y\) is zero, since
that kernel has zero integral homology. Thus it descends to \(Y\).
The two coefficients \(q(b)\), \(q(w+b)\) lift through the actual
surjection \(S_{57}\to\mathcal L\). Conjugating the old images of
\(s,r\) by those two source elements realizes (D.65.9) literally.
Their small and \(\mathcal W_h\) images are zero, and the same two
elements serve at every finite homology level.

The compact change is even, kills all other compact summands, and
has zero small image. It also kills all cusp homology. The full
homology map consequently has the old invertible diagonal
reductions in its cusp/compact filtration. It is an integral
automorphism. The full affine change can have an odd cusp component;
the proof does not replace it by a doubled affine map.

\emph{The integral orientation identity and the actual error.}
Choose actual full-source lifts of the current two conjugators.
Their cusp defect has zero integral abelianization, because the
constructed action fixes \(t\) on the full homology source. It has
an integral class \(z^{\rm int}\) in the class-two commutator lattice.
Apply the marked compact surface boundary word before reduction.
For each cusp, \(t^{36}\) contributes the sum of the 36 conjugates
of this error. Summing over the literal cusp cosets gives the full
group norm. Changes of the cusp conjugator contribute wedges with
cusp homology and vanish under compactification. The commutator
product on the other side of the boundary word is the integral
orientation bivector. Therefore
\[
          N_G(z^{\rm int})_{\rm compact}
             =T\Omega^{\rm int}\bar T^{\,t}-\Omega^{\rm int}.
 \tag{D.65.22}
\]
This is the D.55 proof over the integral lattice, with its original
oriented cusp convention. It uses no interchange of a finite deck
action: the constructed automorphism is identity on \(G\), and
therefore preserves the literal character and pair-coset markings.
The actual conjugator framings determining its compact reductions
are exactly (D.65.13)--(D.65.14).

On the selected induced self block, take the sum of the transported
coset coordinates. Under this actual coinvariant map the left side
of (D.65.22) becomes \(3^{97}N_{H_a}\). This is the full translation
norm factor, including the odd kernel in (D.65.2). The reduction of
the actual projected orientation is the D.64
\(\Omega_a=\overline{(H_a^{\rm int})^{-1}}\); all odd covering and
Fourier factors reduce to one. The modulo-four norm used above is
on this actual integral exterior lattice.

Here is the integral comparison with the smaller graph used in the
certificate. The kernel of \(G\to G^\dagger\) has odd order
\(d=3^{93}\). On each selected character component the actual
pushforward of cover homology is an isomorphism, with inverse
\(d^{-1}\operatorname{tr}\). Indeed pushforward followed by transfer
is the sum of these deck transformations, and transfer followed by
pushforward is multiplication by \(d\); that kernel acts trivially
on the selected component. Compactification commutes with both
maps. Adjunction of the compact intersection pairing gives
\[
 H_G=d^{-1}H_{\dagger},\qquad
 \Omega_G=d\,\Omega_{\dagger}.
 \tag{D.65.23}
\]
Consequently the selected version of (D.65.22), in the literal
certificate coordinates, is
\(N_{H_a}(z)=3^{-4}(T\Omega_{\dagger}\bar T^{\,t}
-\Omega_{\dagger})\). This retains the full integral odd factor.
Here \(\Omega_{\dagger}\) is the actual integral bivector transported
by the specified Fourier map; only its already verified reduction
is needed. Dividing the displayed norm identity by two and reducing
gives the comparison used in (D.65.8).

Divide (D.65.22) by two on this block and reduce modulo two. The
division on the left is defined on the whole lattice by (D.65.6).
For the actual current cusp error write its result as \(n_a(q_az)\).
Replacing the compact homology map by
\(T+\operatorname{Tr}_{H_a}(A)\) changes that result by exactly
(D.65.8). The product of the two trace changes is divisible by four
and disappears after this particular division and reduction. Only
the proved reduction \(D_a\) and the actual reduced orientation are
used; no unspecified integral compact action is substituted.

Apply the retained rank-six right inverse to this actual value
\(n_a(q_az)\), for \(a=2,3\). Lift its ordinary coefficient matrices
to the saturated integral kernel of \(\phi_a\), and use
(D.65.18)--(D.65.21). This defines the actual coefficients in
(D.65.9), and makes both divided-norm errors zero. The two character
supports are disjoint, so their corrections commute at the displayed
linearized level and have no cross term in either selected self
block.

\emph{Retaining the earlier source and killing the full two blocks.}
The compact reductions on every summand are unchanged by the new
correction. The full compact action still has fixed left rows on
the extra projective summands, is identity there, and is block
diagonal on the actual orthogonal decomposition. Thus every earlier
injective-norm argument in D.59--D.64 applies with the same maps.
The earlier self-block corrections may be reapplied in their actual
extra-projective lifts. Their error has an extra-projective factor
or a cusp factor, and hence has zero image in either new compact
self block. Their cusp words give exactly (D.65.10), including all
36 factors and the actual wrap term. They act trivially on integral
homology, so the divided orientation values just computed remain
zero. Central augmentation corrections also leave those values
unchanged, since the divided norm kills all augmentation relations.
This produces a preliminary source lift whose error lies in
\(\mathcal J_{63}\), with \(n_a(q_az)=0\) for both new blocks.

By the exact kernel calculation (D.65.7), each remaining
coinvariant class belongs to the four-dimensional actual
source-involution image. Use the four recorded source vectors in
each orbit and lift them through \(S_{57}\to\mathcal L\).
The literal word is
\(u\widetilde s(u)\), with the current source involution
\(\widetilde s\). Its error is the actual commutator
\([u,\widetilde s(u)]\), giving the D.64 map
\([e_\Delta(v\wedge sv)]_G\). Its \(Y\)-image is fixed by the
current involution, so it leaves (D.65.10) unchanged. These eight
words kill the entire two ten-dimensional coinvariant errors.
Their divided norms are zero by the checked equality of kernels.

After these words, remove any reintroduced moving component in
\((1-e_\Delta)\mathcal K\). This is done after the words, using the
onto actual two-conjugator augmentation map from D.58 and D.64.
Its coefficients lift to \(R_{50}\cap L_{*,2}\subset S_{57}\), so
they preserve \(Y\), the literal small conjugators, and
\(\mathcal W_h\). Their image in each invariant coinvariant test
is zero. Finally apply the same actual augmentation map to the
remaining representatives in
\(C_1\oplus C_2\oplus C_3\oplus\nu_4(C_4)\).
Joint surjectivity of the four maps from D.64 supplies the source
coefficients. The final error lies in \(\mathcal J_{65}\).

The defining kernel is invariant: the compact projections kill all
cusp terms, the old fixed left rows are retained, the actual
\(M_a\)-reductions are unchanged and invertible, and their small
squares commute. The resulting finite class-two map is an
automorphism by its invertible Frattini map. The same two source
conjugators give the integral homology automorphism at every level,
and retain the prescribed \(\mathcal W_h\) automorphism. The actual
marked fibre-product description of D.58 therefore gives an
automorphism of \(X/\mathcal J_{65}\). This proves (D.65.3), with
the literal comparison (D.65.9)--(D.65.10).

\emph{The larger fixed-homology centralizer comparison.}
We give the promised exact maps behind (D.65.11). Restrict
(D.65.19) to \(\langle s\rangle\). Edge modules and order-three
vertex modules are free. The fixed order-two vertices are indexed
by \(C_G(s)/\langle s\rangle\). The two exact chain sequences
identify
\[
 \widehat H^0(\langle s\rangle,U)=0,\qquad
 \widehat H^{-1}(\langle s\rangle,U)
   =\ker\left(\mathbb F_2[C_G(s)/\langle s\rangle]
                                    \longrightarrow\mathbb F_2\right).
 \tag{D.65.24}
\]
The map to the small module is the literal map of fixed vertices.
Centralizers lift onto across the elementary-abelian odd kernel of
\(G\to G^\dagger\). Explicitly, a lifted centralizing element has
\(gsg^{-1}=ks\), with \(sk=-k\) in that kernel. Replacing \(g\)
by \((-k/2)g\) makes it centralize \(s\). At the remaining finite step,
\(C_Q(s)\to C_{Q_*}(s_*)\) is onto, with orders 16 and 8.
The common fixed translation space has nine elements. At the
\(G^\dagger\to G_*\) stage the actual fixed-vertex map has 72 source
vertices, 36 target vertices, and fibre size two. Its reduced
augmentation map is onto of rank
35. Every coset and an actual section are retained in
\path{d_source_centralizer_norm.py}. The Tate exact sequence of
\(0\to\ker(U\to U_*)\to U\to U_*\to0\) now gives zero
\(\widehat H^0\) for the kernel. The actual odd-character projector
commutes with the kernel map; its direct summand \(\mathcal L\)
has the same vanishing, proving the first equality
in (D.65.11).

For the cusp Sylow subgroup of order four, both vertex permutation
modules are free. The unique involution has the actual \(Q\)-state
314; it is not conjugate to the source state 1, in either \(Q\) or
\(Q_*\). The full conjugacy classes are checked. Hence it fixes no
order-two vertex; order-three vertex freeness follows from the
orders. The two chain sequences identify the even Tate group of
\(U\) with \(\mathbb Z/4\) and the odd group with zero. The map to
\(U_*\) is the identity on the terminal augmentation lattice, so
it is an isomorphism on these groups. The kernel is consequently
cohomologically trivial for this cyclic subgroup. Averaging over
the odd part of \(\langle t\rangle\) proves both cusp equalities
in (D.65.11). The recorded affine cusp has order twelve in
\(G^\dagger\); its sixth power is state 314 with translation
\((0,2,1,0)\). The actual source cusp retains order 36 throughout.

An element of the kernel of the literal small comparison in \(Y\)
has homology in \(\mathcal L\). If it centralizes the current
involution or cusp, its homology is fixed by the indicated cyclic
action. Equation (D.65.11) gives precisely a norm vector. Its
actual norm word has that homology; the remaining central difference
changes the error only by an augmentation relation. Thus the
fixed-homology route has exactly the source maps already tested.
The new construction instead has the explicit comparison
(D.65.9)--(D.65.10).

\emph{The remaining kernel and the execution record.}
The four maps of D.64 are jointly onto. Replacing the retractions
\(p_2,p_3\) by the full second and third coordinates therefore
gives (D.65.4), with
\(2(2880-4)=5752\) and \(2(10-4)=12\).
The total retained quotient has dimension
\(1920+2880+2880+2=7682\) and coinvariant dimension
\(4+10+10+2=26\). The specified residual dimension is
\(282230-5752=276478\). Joint surjectivity also maps
\(\mathcal J_{65}\) onto \(\ker\nu_4\), whose dimension is 1918.
This proves the nonzero source-kernel assertion through its actual
preimage in (D.65.12).

Eight complete programs retain all finite columns, actual framing
words, integral graph maps, inverse matrices, norm relations,
linearized equations and correction sections. The separate
permutation comparison and real-exterior verification are included.
An initial integral run completed the first mathematical block but
failed while serializing a variable that had been reused as an
index; its source and actual output are preserved. The corrected
program uses a distinct section variable. No failed run is counted
as a completed certificate. The final copied-source replay and
all input hashes are recorded separately.
\end{proof}
